\documentclass[preview]{standalone}

\usepackage{amsmath}
\usepackage{amssymb}
\usepackage{stellar}
\usepackage{definitions}

\begin{document}

\id{psicologia-teorie-intelligenza}
\genpage

\section{Teorie dell'intelligenza}

\begin{snippetdefinition}{intelligenza-definition}{Intelligenza}
    L'\emph{intelligenza} è un'abilità mentale generale che comprende
    ragionamento, pianificazione, soluzione di problemi, pensiero astratto,
    comprensione di idee complesse, apprendimento rapido e apprendimento
    dall'esperienza.
\end{snippetdefinition}

\begin{snippet}{intelligenza-costrutto}
    L'intelligenza è un costrutto teorico ampio e difficile da caratterizzare.
    Le sue definizioni oscillano tra capacità di ragionamento astratto,
    adattamento all'ambiente, apprendimento e gestione efficace di situazioni
    nuove o complesse.
\end{snippet}

\begin{snippetdefinition}{teorie-esplicite-intelligenza-definition}{Teorie esplicite dell'intelligenza}
    Le \emph{teorie esplicite dell'intelligenza} sono teorie formulate da
    esperti e fondate sulle risposte di gruppi di soggetti a compiti considerati
    idonei a misurare il comportamento intelligente.
\end{snippetdefinition}

\begin{snippetdefinition}{teorie-implicite-intelligenza-definition}{Teorie implicite dell'intelligenza}
    Le \emph{teorie implicite dell'intelligenza} riguardano le idee che le
    persone hanno sulla propria intelligenza e su quella altrui. Possono
    influenzare la costruzione delle teorie esplicite e mettere in luce
    differenze culturali nelle concezioni dell'intelligenza.
\end{snippetdefinition}

\begin{snippetdefinition}{analisi-fattoriale-definition}{Analisi fattoriale}
    L'\emph{analisi fattoriale} è un insieme di tecniche
    matematico-statistiche che permette di sintetizzare le informazioni espresse
    da molte variabili non indipendenti in un numero inferiore di variabili
    ipotetiche soggiacenti.
\end{snippetdefinition}

\begin{snippetexample}{capacita-induttiva-example}{Capacità induttiva}
    Se i componenti di un campione ottengono punteggi simili in test di
    ragionamento induttivo, come classificare, costruire serie logiche o
    rispondere a quesiti analogici, l'analisi fattoriale può individuare una
    capacità induttiva comune alla base delle prestazioni.
\end{snippetexample}

\includesnpt[width=52\%,src=/snippet/static-psychology/factor-analysis-intelligence.jpg]{centered-img}

\begin{snippetdefinition}{teoria-bifattoriale-spearman-definition}{Teoria bifattoriale di Spearman}
    La \emph{teoria bifattoriale di Spearman} sostiene che l'intelligenza sia
    espressione di un fattore generale \(g\), comune alle prestazioni
    intellettive, e di fattori specifici \(s\), legati alle singole abilità o
    prove.
\end{snippetdefinition}

\begin{snippetdefinition}{abilita-mentali-primarie-definition}{Abilità mentali primarie}
    Le \emph{abilità mentali primarie}, o \emph{PMA}, sono le dimensioni
    intellettive individuate da Thurstone: comprensione verbale, fluidità
    verbale, abilità numerica, relazioni spaziali, memoria associativa, velocità
    percettiva e induzione.
\end{snippetdefinition}

\begin{snippetdefinition}{intelligenza-fluida-definition}{Intelligenza fluida}
    L'\emph{intelligenza fluida}, o \(Gf\), è l'abilità di elaborare
    informazioni e ragionare in modo astratto. È connessa ad aspetti biologici
    ed è quindi maggiormente soggetta a deterioramento.
\end{snippetdefinition}

\begin{snippetdefinition}{intelligenza-cristallizzata-definition}{Intelligenza cristallizzata}
    L'\emph{intelligenza cristallizzata}, o \(Gc\), è costituita dai contenuti
    appresi nel corso dell'esperienza, come vocabolario, comprensione verbale e
    informazioni generali.
\end{snippetdefinition}

\begin{snippet}{cattell-horn-ciclo-vita}
    Secondo Cattell e Horn, le abilità cristallizzate possono aumentare nel
    corso del ciclo di vita, mentre le abilità fluide crescono negli anni
    giovanili e tendono a diminuire in età avanzata.
\end{snippet}

\includesnpt[width=52\%,src=/snippet/static-psychology/fluid-crystallized-intelligence-lifespan.jpg]{centered-img}

\begin{snippetdefinition}{teoria-triarchica-intelligenza-definition}{Teoria triarchica dell'intelligenza}
    La \emph{teoria triarchica dell'intelligenza} di Sternberg descrive
    l'intelligenza attraverso tre dimensioni: componenziale, esperienziale e
    contestuale.
\end{snippetdefinition}

\begin{snippet}{sternberg-subteorie}
    La dimensione componenziale riguarda i meccanismi mentali interni che
    sostengono il comportamento intelligente. La dimensione esperienziale
    riguarda l'uso dell'esperienza, la gestione della novità e
    l'automatizzazione. La dimensione contestuale riguarda adattamento, selezione
    e modellazione del contesto socioculturale.
\end{snippet}

\begin{snippet}{sternberg-intelligenze}
    Sternberg distingue anche tre generi di attività intelligente:
    \begin{itemize}
        \item \emph{analitica}, legata ad analisi, giudizio e valutazione;
        \item \emph{creativa}, legata a invenzione, scoperta e immaginazione;
        \item \emph{pratica}, legata all'uso di strumenti e all'attuazione di
        progetti.
    \end{itemize}
\end{snippet}

\includesnpt[width=58\%,src=/snippet/static-psychology/sternberg-triarchic-intelligence.jpg]{centered-img}

\begin{snippetdefinition}{teoria-intelligenze-multiple-definition}{Teoria delle intelligenze multiple}
    La \emph{teoria delle intelligenze multiple} di Gardner sostiene che non
    esista una sola intelligenza generale, ma una pluralità di intelligenze
    relativamente indipendenti, combinate in modo peculiare in ogni individuo.
\end{snippetdefinition}

\begin{snippet}{gardner-intelligenze}
    Nella revisione del 1999, Gardner individua nove tipi di intelligenza,
    ciascuno relativamente indipendente dagli altri e ancorato a un proprio
    substrato neurologico. Ogni persona possiede una combinazione specifica di
    intelligenze, più o meno efficace a seconda dei contesti.
\end{snippet}

\includesnpt[width=55\%,src=/snippet/static-psychology/gardner-multiple-intelligences.jpg]{centered-img}

\begin{snippetdefinition}{intelligenza-emotiva-definition}{Intelligenza emotiva}
    L'\emph{intelligenza emotiva} è l'insieme di abilità legate alla gestione
    delle emozioni proprie e altrui, come controllo degli impulsi,
    automotivazione, empatia e competenza sociale.
\end{snippetdefinition}

\begin{snippet}{intelligenza-emotiva-valutazione}
    Nel costrutto di intelligenza emotiva si intrecciano aspetti cognitivi e
    affettivi. Resta quindi aperta la questione di quanto un comportamento
    emotivamente intelligente rifletta una capacità intellettiva specifica e
    quanto dipenda da personalità, stabilità emotiva e adattamento sociale.
\end{snippet}

\end{document}
